\subsection{LLM 기반 에이전트 워크플로우 설계 기법}
최근 대규모 언어 모델(LLM)의 활용 패러다임은 단순한 단일 프롬프트 기반의 질의응답에서 벗어나, 복잡한 목표를 달성하기 위해 일련의 단계적 절차를 수행하는 에이전트 워크플로우(Agentic Workflow) 설계로 진화하고 있다. 초기 연구들이 주로 Chain-of-Thought(CoT)와 같은 추론 경로의 명시화를 통해 모델의 성능을 높이는 데 집중했다면, 최근의 경향은 특정 역할(Role)을 부여받은 다수의 에이전트가 협력하는 멀티 에이전트 시스템(Multi-Agent System, MAS)으로 확장되는 추세이다. 특히 문서 생성 작업에서는 설계자(Planner), 집필자(Writer), 검수자(Reviewer)와 같이 세분화된 역할을 정의하고, 이들 간의 상호작용을 통해 결과물의 품질을 점진적으로 개선하는 반복적 정제(Iterative Refinement) 기법이 핵심적으로 다루어지고 있다. 이러한 워크플로우 설계의 핵심은 각 에이전트의 전문성을 극대화하면서도, 전체 프로세스의 일관성을 유지하기 위한 제어 메커니즘을 어떻게 설계하느냐에 있다.

\subsection{자율적 에이전트 구조와 제어의 한계}
에이전트 시스템의 자율성을 극대화하려는 시도로 AutoGPT와 BabyAGI와 같은 자율적 에이전트 구조가 제안되었다. 이러한 시스템은 목표 달성을 위해 작업 목록을 스스로 생성하고 우선순위를 조정하는 재귀적 루프(Recursive Loop) 구조를 통해, 사용자가 제시한 최종 목표를 달성하기 위한 하위 과업(Sub-task) 생성 및 실행 순서 결정을 동적으로 수행한다. 그러나 이러한 완전 자율형 구조는 엄격한 형식이 요구되는 산업 현장의 정형 문서 생성 작업에 적용할 때 명확한 한계를 드러낸다. 첫째, LLM의 추론 오류로 인해 불필요한 하위 과업을 무한히 생성하거나 목표에서 벗어난 경로로 진입하는 '탈주(Drift)' 현상이 발생하여 제어 가능성이 낮다. 둘째, 동적 라우팅 과정에서 발생하는 반복적인 자기 성찰(Self-reflection) 단계로 인해 토큰 소비가 급증하며 자원 효율성이 저하된다. 따라서 최근의 연구들은 완전한 자율성보다는 적절한 제약 조건(Constraint) 하에서 LLM의 유연성을 활용하는 '제어된 자율성(Controlled Autonomy)' 방향으로 선회하고 있으며, 이는 본 연구에서 제안하는 오케스트레이션 구조의 설계 배경이 된다.

\subsection{정형 문서 생성 자동화 및 본 연구의 학술적 위치}
기존의 정형 보고서와 같은 전문 문서 생성 자동화 연구들은 주로 템플릿 기반의 생성 방식이나 검색 증강 생성(RAG)을 통한 지식 보강에 집중해 왔으며, 본 연구는 이러한 접근 방식을 넘어 오케스트레이션 구조의 성능 비교를 통해 새로운 시각을 제시한다. 산업 현장의 보고서는 단순한 정보의 나열을 넘어 논리적 구조의 완결성과 엄격한 형식 준수가 필수적이므로, 본 연구는 오케스트레이션 구조를 세 가지 관점에서 정의하여 분석을 수행한다.

먼저, 본 연구의 \textbf{Code 기반 구조}는 전통적인 소프트웨어 공학의 결정론적 워크플로우(Deterministic Workflow) 및 Directed Acyclic Graph(DAG) 기반의 파이프라인 설계 방식에 위치한다. 이 구조는 결정론적 특성으로 인해 높은 예측 가능성과 실행 효율성을 제공한다. 실제 실험 결과, Code 기반 구조는 평균 수행 시간 158.61초, 토큰 소비량 21,440개, 최고 점수(\textit{best\_score}) 82점을 기록하며, 정형화된 프로세스가 확립된 작업에서의 실질적인 효율성을 입증하였다.

둘째, \textbf{LLM 기반 구조}는 앞서 언급한 자율적 에이전트의 동적 라우팅(Dynamic Routing) 개념을 계승한다. 이는 작업의 복잡도나 중간 결과물의 품질에 따라 경로를 유연하게 변경하는 방식으로, 창의적 글쓰기나 비정형 작업에서 높은 잠재력을 가진다.

마지막으로, \textbf{Hybrid 구조}는 상위 수준의 전략적 제어는 코드(Code)로 수행하고 하위 수준의 전술적 실행은 LLM의 판단에 맡기는 계층적 제어(Hierarchical Control) 구조를 취한다. 이는 결정론적 구조의 안정성과 동적 구조의 유연성을 결합한 형태로, 고도의 정밀함과 유연함이 동시에 요구되는 산업용 전문 문서 생성의 최적 대안으로서 학술적 위치를 점한다. 본 연구는 이 세 가지 구조를 동일한 통제 변인 하에서 정량적으로 비교함으로써, 작업의 성격에 따른 최적의 오케스트레이션 선택 기준을 제시한다는 점에서 기존의 단일 구조 제안 연구들과 차별성을 가진다.